\chapter{First-Order Deduction}
\label{ch:fol-deduction}

\section{Natural-deduction rules for quantifiers and equality}
\label{sec:fol-deduction-01}

The calculus \(\ND_{\FOL}\) contains every rule of \(\ND_{\CPL}\) and the
following rules.  Every displayed substitution is the capture-avoiding
substitution of \cref{ch:fol-syntax-semantics}.
\[
\begin{array}{c@{\qquad}c}
\Rule[\(\forall E\)]{\forall x\,\varphi}{\varphi[t/x]}
&
\Rule[\(\exists I\)]{\varphi[t/x]}{\exists x\,\varphi}
\end{array}
\]
\[
\Rule[\(\forall I\)]{\varphi}{\forall x\,\varphi}
\side{provided \(x\) is not free in any open assumption}.
\]
\[
\Rule[\(\exists E\)]{
 \exists x\,\varphi\qquad
 \begin{array}{c}[\varphi[y/x]]^i\\ \vdots\\ \psi\end{array}}
 {\psi}
\side{\(y\) is absent from \(\psi\), from every other open assumption,
and from \(\exists x\,\varphi\)}.
\]
A variable subject to a freshness side condition in a rule instance is the
\emph{eigenvariable} of that instance; in \(\exists E\) the eigenvariable
is \(y\), and in \(\forall I\) it is the generalized variable \(x\).  The
rule \(\exists E\) discharges
all assumptions marked \(i\).  Its freshness condition prevents the temporary
witness from appearing in the conclusion or surviving assumptions.

Because formulas are strings rather than alpha-equivalence classes, the
calculus also contains the explicit conversion rule
\[
  \Rule[\(\alpha\)]{\varphi}{\psi}
  \side{provided \(\varphi\equiv_\alpha\psi\)}.
\]
The rule's premise and conclusion are formulas, and its side condition is the
syntactic alpha-equivalence relation of
\cref{sec:fol-syntax-semantics-04}.  It changes only bound-variable names, so
it is neither identity elimination nor replacement by an arbitrary
logically equivalent formula.  The rule is required below when a fresh
eigenvariable of \(\forall I\) produces an alpha-equivalent, but not
literally identical, quantified formula.

Equality is governed by identity introduction and substitutivity:
\[
\begin{array}{c@{\qquad}c}
\Rule[\(=I\)]{}{t=t}
&
\Rule[\(=E\)]{r=t\qquad\varphi[r/x]}{\varphi[t/x]}.
\end{array}
\]
The formula \(\varphi\) may contain zero or more free occurrences of \(x\).
The rule uses actual identity in the semantics; it does not merely postulate
that an arbitrary relation is reflexive.

\begin{definition}[First-order derivability]
For \(\Gamma\subseteq\Fm(\Lang)\), write
\(\Gamma\derives_{\ND_{\FOL}}\varphi\) iff there is a
\(\ND_{\FOL}\)-derivation of \(\varphi\) whose open assumptions all belong to
\(\Gamma\).
\end{definition}

\begin{lemma}[Invariance under reduct]
\label{lem:reduct-invariance}
Let \(\Lang'\) extend \(\Lang\), let \(\Meta'\) be a
\(\Lang'\)-structure, and let \(s\) be a variable assignment in its domain.
For every \(\Lang\)-term \(t\) and \(\Lang\)-formula \(\varphi\),
\[
 \sem{t}^{\Meta'}_s=\sem{t}^{\Meta'\restr\Lang}_s,
 \qquad
 \Meta',s\entails\varphi
 \quad\Longleftrightarrow\quad
 \Meta'\restr\Lang,s\entails\varphi.
\]
\end{lemma}
\begin{proof}
Induct on terms for the first equality.  Every function symbol in the term has
the same interpretation in the expansion and reduct.  Then induct on formulas.
The atomic cases use the term result and equality of the old relation-symbol
interpretations.  The clauses for \(\bottom\) and the connectives are
immediate from the induction hypotheses.  In a quantifier case the two
structures have the same domain, so they range over the same elements; apply
the induction hypothesis under each updated assignment.
\end{proof}

\begin{example}[Universal reasoning]
The premises
\(\forall x(P(x)\to Q(x))\) and \(\forall xP(x)\) derive
\(\forall xQ(x)\):
\[
\begin{array}{rll}
1.&\forall x(P(x)\to Q(x))&\text{premise}\cr
2.&\forall xP(x)&\text{premise}\cr
3.&P(y)\to Q(y)&\forall E,1\cr
4.&P(y)&\forall E,2\cr
5.&Q(y)&\to E,3,4\cr
6.&\forall yQ(y)&\forall I,5\cr
7.&\forall xQ(x)&\alpha,6.
\end{array}
\]
The variable \(y\) is not free in either premise, so the side condition at
line 6 is met.  Line 7 applies the \(\alpha\) rule, because
\(\forall yQ(y)\) and \(\forall xQ(x)\) are alpha-equivalent but not
identical strings.
\end{example}

\begin{example}[Existential elimination]
From \(\exists x(P(x)\land Q(x))\) derive \(\exists xP(x)\):
\[
\begin{array}{rll}
1.&\exists x(P(x)\land Q(x))&\text{premise}\cr
2.&\quad[P(y)\land Q(y)]^1&\text{assumption, \(y\) fresh}\cr
3.&\quad P(y)&\land E_1,2\cr
4.&\quad\exists xP(x)&\exists I,3\cr
5.&\exists xP(x)&\exists E,1,2\text{--}4.
\end{array}
\]
The eigenvariable \(y\) does not occur in the conclusion.
\end{example}

\begin{example}[Symmetry of equality]
The rule \(=E\) derives \(t=r\) from \(r=t\).  Choose a variable \(x\)
with \(x\notin\operatorname{Var}(r)\); this is possible because \(\Var\)
is infinite and \(r\) is a finite string.  From \(=I\) obtain
\(r=r\).  In the formula \(\varphi\defeq x=r\), the instance
\(\varphi[r/x]\) is \(r=r\) and the instance \(\varphi[t/x]\) is
\(t=r\).  Apply \(=E\) to \(r=t\) and \(r=r\).
\end{example}

